% Options for packages loaded elsewhere
\PassOptionsToPackage{unicode}{hyperref}
\PassOptionsToPackage{hyphens}{url}
%
\documentclass[
  11pt,
]{article}
\usepackage{lmodern}
\usepackage{amssymb,amsmath}
\usepackage{ifxetex,ifluatex}
\ifnum 0\ifxetex 1\fi\ifluatex 1\fi=0 % if pdftex
  \usepackage[T1]{fontenc}
  \usepackage[utf8]{inputenc}
  \usepackage{textcomp} % provide euro and other symbols
\else % if luatex or xetex
  \usepackage{unicode-math}
  \defaultfontfeatures{Scale=MatchLowercase}
  \defaultfontfeatures[\rmfamily]{Ligatures=TeX,Scale=1}
\fi
% Use upquote if available, for straight quotes in verbatim environments
\IfFileExists{upquote.sty}{\usepackage{upquote}}{}
\IfFileExists{microtype.sty}{% use microtype if available
  \usepackage[]{microtype}
  \UseMicrotypeSet[protrusion]{basicmath} % disable protrusion for tt fonts
}{}
\makeatletter
\@ifundefined{KOMAClassName}{% if non-KOMA class
  \IfFileExists{parskip.sty}{%
    \usepackage{parskip}
  }{% else
    \setlength{\parindent}{0pt}
    \setlength{\parskip}{6pt plus 2pt minus 1pt}}
}{% if KOMA class
  \KOMAoptions{parskip=half}}
\makeatother
\usepackage{xcolor}
\IfFileExists{xurl.sty}{\usepackage{xurl}}{} % add URL line breaks if available
\IfFileExists{bookmark.sty}{\usepackage{bookmark}}{\usepackage{hyperref}}
\hypersetup{
  hidelinks,
  pdfcreator={LaTeX via pandoc}}
\urlstyle{same} % disable monospaced font for URLs
\usepackage[margin=1in]{geometry}
\setlength{\emergencystretch}{3em} % prevent overfull lines
\providecommand{\tightlist}{%
  \setlength{\itemsep}{0pt}\setlength{\parskip}{0pt}}
\setcounter{secnumdepth}{-\maxdimen} % remove section numbering
\usepackage{esint}
\setmathfont{Asana Math}
\setlength{\emergencystretch}{3em}

\author{}
\date{}

\begin{document}

\hypertarget{topological-local-regularity-and-a-diameter-free-einstein-repair-criterion}{%
\section{Topological local regularity and a diameter-free Einstein
repair
criterion}\label{topological-local-regularity-and-a-diameter-free-einstein-repair-criterion}}

24 September 2026. This note separates a local geometric extension from
the classical global elliptic solver it can support. It does not infer a
uniform inverse from a curvature bound or from absence of a known
obstruction.

\hypertarget{the-geometric-statement-needs-neither-diameter-nor-a-weyl-bound}{%
\subsection{1. The geometric statement needs neither diameter nor a Weyl
bound}\label{the-geometric-statement-needs-neither-diameter-nor-a-weyl-bound}}

Write \(V_{-K/3}(r)\) for the four-dimensional model-ball volume,
\(\omega_4=\pi^2/2\), and \(\theta=v/V_{-K/3}(1)\). For a nonempty class
below \(0<\theta\le1\).

\textbf{Theorem 1 (diameter-free local regularity).} Fix \(K,v>0\) and
\(0<\alpha<1\). Closed connected oriented smooth four-manifolds with \[
 \|\operatorname{Ric}^{\sharp}\|_{\infty,\mathrm{op}}\le K,
 \qquad \inf_{x\in M}\operatorname{Vol}B_1(x)\ge v,                \tag{2}
\] and \(\pi_2(M)=0\) have a uniform positive \(C^{1,\alpha}\) harmonic
radius. Its bound depends only on \(K,v,\alpha\), not on diameter, total
volume or topology. No Weyl bound or condition on \(\pi_1(M)\) is
required. No finite number of smooth types or global compactness is
asserted without a diameter bound.

\textbf{Proof.} Relative Bishop comparison gives \[
 \operatorname{Vol}B_r(x)\ge
 v\frac{V_{-K/3}(r)}{V_{-K/3}(1)}\ge\theta\omega_4r^4,
 \qquad 0<r\le1.                                             \tag{4}
\] The local Jiang--Naber curvature estimate, after a fixed rescaling,
gives \[
 \int_{B_a(x)}|\operatorname{Rm}|^2\le C(K,v)                  \tag{5}
\] for a fixed \(a=a(K)>0\). Only a ball around the eventual blowup
centre is used; a global covering and a total-energy bound are
unnecessary.

Suppose the minimum harmonic radius \(\rho_j\) tends to zero, at
\(x_j\). Rescale by \(\rho_j^{-2}\). Every rescaled harmonic radius is
at least one, Ricci tends uniformly to zero, and (4) gives Euclidean
lower volume growth on every fixed rescaled ball. The
harmonic-coordinate argument in
\href{20260923_compact_realization.md}{the compact realization proof},
section 4, therefore gives a complete nonflat smooth Ricci-flat pointed
limit \(X\). In particular, the harmonic metric equation gives strong
local \(W^{2,p}\) convergence, for each finite \(p\), without
differentiating the Ricci error. The harmonic-radius normalization
excludes a flat limit.

For each fixed \(R\), the rescaled radius-\(R\) ball lies inside the
original \(B_a(x_j)\) for large \(j\). Passing (5) to the limit and
exhausting in \(R\) gives finite \textbf{total} curvature energy on
\(X\), with bound \(C(K,v)\). The lower volume ratio is at least
\(\theta\). The Ricci-flat asymptotic theorem then gives an ALE end of
positive order. Optimal decay order is not needed.

A compact core \(B\) of \(X\) embeds in \(M_j\). The strengthened
topological input is
\href{https://arxiv.org/pdf/2305.03154v3}{Bamler--Chen, Lemma 4.35}: no
nontrivial Ricci-flat ALE core embeds in a smooth four-manifold with
integer \(H_2=0\) and torsion-free integer \(H_1\). The proof cites
Zhang's Theorem 6.1 and the earlier Anderson argument; it does not
assume half-flatness.

For completeness, applying that lemma here does not require the original
fundamental group to be torsion-free. A Ricci-flat ALE space has finite
fundamental group. A connected component of the preimage of \(B\) in
\(\widetilde M_j\) is therefore a finite cover \(\widehat B\) of \(B\),
embedded in \(\widetilde M_j\). Its interior carries the lifted complete
nonflat Ricci-flat ALE metric: the cover is finite and the end remains
an ALE end. Hurewicz and \(\pi_2(M_j)=0\) give
\(H_1(\widetilde M_j;\mathbb Z)=H_2(\widetilde M_j;\mathbb Z)=0\). Lemma
4.35 now excludes \(\widehat B\), the required contradiction. This
proves the harmonic-radius bound. \(\square\)

The finite-cover use of this lemma is developed separately in
\href{20260924_ale_boundary_compatibility.md}{the topological
continuation}.

The external curvature input in (5) is
\href{https://arxiv.org/html/1605.05583}{Jiang--Naber, Theorem 1.6}. It
controls local full curvature energy from bounded Ricci and noncollapse;
it is not an Einstein regularity theorem imposed on the input metrics.
The local use of that theorem is the reason diameter disappears from
this part of the argument.

This stronger classical topological input supersedes the need for chiral
concentration control in this local regularity theorem. The older
half-flat argument remains valid but is not the strongest available one.

\hypertarget{what-replaces-compactness-in-a-global-existence-argument}{%
\subsection{2. What replaces compactness in a global existence
argument}\label{what-replaces-compactness-in-a-global-existence-argument}}

Theorem 1 supplies local geometry. A global inverse is a different
input. We give a sufficient condition that yields an exact metric,
rather than another necessary condition on a hypothetical output.

For a smooth background metric \(g\), let \[
 B_gu=\delta_gu+\tfrac12d\operatorname{tr}_gu,
 \qquad
 \Phi_g(h)=\operatorname{Ric}_h+h+\delta_h^*(B_gh).               \tag{9}
\] Here divergence is minus covariant contraction and \(\delta_h^*\) is
its formal adjoint on one-forms. The derivative at \(g\) is the
self-adjoint elliptic operator \[
 L_g=\tfrac12\Delta_{L,g}+I,
 \quad
 \Delta_{L,g}u=\nabla^*\nabla u+\operatorname{Ric}\circ u
             +u\circ\operatorname{Ric}-2\mathring R u.          \tag{10}
\] An \textbf{absolute spectral gap} means \[
 \|L_gu\|_{L^2}\ge\lambda\|u\|_{L^2}
       \quad\text{for every smooth symmetric two-tensor }u.     \tag{11}
\] This permits negative eigenvalues but excludes both a kernel and
small eigenvalues of either sign. It is not a scalar Laplacian gap.

\textbf{Theorem 2 (topological local geometry plus a global inverse).}
Fix \(v,K_1,\lambda>0\) and \(0<\alpha<1\). There exist
\(\varepsilon,C>0\), depending only on these data, such that the
following holds. Let a closed connected oriented smooth four-manifold
with \(\pi_2(M)=0\) carry a smooth metric satisfying \[
 \inf_x\operatorname{Vol}B_1(x)\ge v,\qquad
 \|\nabla\operatorname{Ric}\|_\infty\le K_1,
 \qquad (11).                                                \tag{12}
\] Write \(E=\operatorname{Ric}_g+g\) and
\(\delta=\|g^{-1}E\|_{\infty,\mathrm{op}}\). If \[
 \delta^{1-\alpha}+\|E\|_{L^2(M,g)}<\varepsilon,                \tag{13}
\] then there is a smooth metric \(h\) on the same smooth manifold with
\[
 \operatorname{Ric}_h=-h,\qquad
 \|h-g\|_{C^{2,\alpha}}+\|h-g\|_{H^2}
       \le C\bigl(\delta^{1-\alpha}+\|E\|_2\bigr).             \tag{14}
\] The Hölder norm is measured in uniformly controlled small harmonic
charts; an equivalent intrinsic bounded-geometry norm gives the same
statement. There is no diameter or total-volume bound. Neither a Kähler
output nor negative sectional curvature is required. Condition (13) uses
the \textbf{unnormalized} integral residual, not its mean.

\textbf{Proof.} Take \(\varepsilon\) small enough that \(\delta<1/2\).
Then the Ricci bound needed in Theorem 1 holds with \(K=2\). First use
its harmonic radius conclusion with an exponent strictly larger than
\(\alpha\). In these coordinates the bound on
\(\nabla\operatorname{Ric}\), the controlled Christoffel symbols and the
Ricci bound control the coordinate \(C^{0,\alpha}\) norm of Ricci.
Schauder estimates in the harmonic metric equation therefore give
uniform \(C^{2,\alpha}\) bounds for the metric on smaller charts, and a
uniform curvature bound. These local constants have no diameter or
volume dependence.

Use the Banach spaces \[
 X=C^{2,\alpha}\cap H^2,\qquad Y=C^{0,\alpha}\cap L^2,
 \quad \|u\|_X=\|u\|_{C^{2,\alpha}}+\|u\|_{H^2},
 \quad \|f\|_Y=\|f\|_{C^{0,\alpha}}+\|f\|_2.                   \tag{15}
\] For each individual closed manifold, (11) and self-adjoint elliptic
Fredholm theory give a unique inverse and
\(\|L_g^{-1}f\|_2\le\lambda^{-1}\|f\|_2\). Uniform local elliptic
estimates, summed in a bounded-overlap atlas, give \[
 \|u\|_{H^2}\le C(\|L_gu\|_2+\|u\|_2).                      \tag{16}
\] The atlas may have arbitrarily many members; its overlap bound,
rather than its cardinality, enters this inequality. Local Moser
estimates for a connection Laplacian with bounded zero-order
coefficients, followed by Schauder estimates, give \[
 \|u\|_{C^{2,\alpha}}
       \le C(\|L_gu\|_{C^{0,\alpha}}+\|u\|_2).                \tag{17}
\] Thus \[
                         \|L_g^{-1}f\|_X\le C_I\|f\|_Y       \tag{18}
\] uniformly across (12). Local bounded-geometry estimates and the
explicit global hypothesis (11), not compactness of the whole manifold,
give (18).

Expand \[
 \Phi_g(g+u)=E+L_gu+N_g(u).
\] For \(\|u\|_X,\|z\|_X\le\rho\), with a fixed small \(\rho\), the
standard coordinate formula for Ricci and (9) gives \[
 \|N_g(u)-N_g(z)\|_Y
     \le C_N(\|u\|_X+\|z\|_X)\|u-z\|_X.                    \tag{19}
\] For clarity, its second-order terms have the form \(u*\nabla^2u\),
\(\nabla u*\nabla u\), and bounded curvature times \(u*u\), with smooth
coefficients for positive-definite \(g+u\). Hölder product estimates
bound their first norm in (15); their \(L^2\) norm is bounded by one
\(C^{2,\alpha}\) factor and one \(H^2\) factor. There is no
multiplication by \(\operatorname{Vol}(M)^{1/2}\).

Tensor interpolation, using \(\nabla E=\nabla\operatorname{Ric}\),
yields \[
 \|E\|_{C^{0,\alpha}}\le C\delta^{1-\alpha},\qquad \delta\le1,
                                                                  \tag{20}
\] where \(C\) can depend on \(K_1\) and the chart data. On the ball of
radius \(R=2C_I\|E\|_Y\) in \(X\), the map
\(u\mapsto-L_g^{-1}(E+N_g(u))\) is a contraction when \(R\le\rho\) and
\(2C_IC_NR<1\). This gives a positive-definite solution \(h=g+u\) of
\(\Phi_g(h)=0\) and (14).

Smallness in \(C^{2,\alpha}\) makes \(\operatorname{Ric}_h\) negative
definite. Put \(b=B_gh\). Applying the contracted Bianchi operator for
\(h\) gives \[
 0=2B_h\delta_h^*b=\nabla_h^*\nabla_h b-\operatorname{Ric}_h(b).
\] Integration on the closed manifold forces \(b=0\). Hence
\(\operatorname{Ric}_h=-h\). Elliptic bootstrapping of the gauged
equation with smooth background makes \(h\) smooth. \(\square\)

The Bianchi-gauge operator, local Moser estimate and hybrid-norm repair
method are classical. For primary statements see
\href{https://arxiv.org/html/2502.16692v2}{Jäckel, sections 2 and 4--5}
and \href{https://numdam.org/item/10.5802/tsg.367.pdf}{Premoselli,
section 4.1}. The contribution here is the geometric input to this
solver: topological core exclusion and local noncollapse supply
controlled charts without global diameter. The spectral gap remains an
independent hypothesis.

\hypertarget{a-pointwise-input-application-with-localized-error}{%
\subsection{3. A pointwise-input application with localized
error}\label{a-pointwise-input-application-with-localized-error}}

\textbf{Corollary 3.} In Theorem 2, suppose \(E\) vanishes outside a
measurable set \(\Omega\) of volume at most \(V_0\). A sufficiently
small pointwise defect \(\delta<\delta_0(v,K_1,\lambda,V_0,\alpha)\)
then gives an exact negative Einstein metric on the same smooth
manifold, with no bound on the total volume or diameter.

Indeed, in four dimensions \(|E|\le2\delta\), so \[
 \|E\|_2\le2\delta\sqrt{V_0}.                               \tag{21}
\] Take \(\delta_0\) to satisfy (13). Exact Einstein regions of
arbitrarily large volume do not increase the residual bound. Their
possible effect on small global eigenvalues must still be tested through
(11).

This gives a sufficient same-smooth-type realization statement, not a
claim that every original Yau26 input has localized error. The repaired
metric is not forced to respect an input symmetry or complex structure.

\hypertarget{which-information-cannot-be-discarded}{%
\subsection{4. Which information cannot be
discarded}\label{which-information-cannot-be-discarded}}

\begin{enumerate}
\def\labelenumi{\arabic{enumi}.}
\tightlist
\item
  \textbf{Averaging does not provide (13).} If
  \(r=\fint|\operatorname{Ric}+g|^2\), then
  \(\|E\|_2=\sqrt{r\operatorname{Vol}(M)}\). Small \(r\) with growing
  volume is insufficient. For an explicit normalization test, take
  covers \((M_j,h_j)\) of degree \(N_j\to\infty\) of a fixed closed
  hyperbolic four-manifold, scaled to \(\operatorname{Ric}_{h_j}=-h_j\),
  and put \(g_j=(1+t_j)h_j\), \(t_j=N_j^{-1/4}\). Then \[
   \delta_-(g_j)=\frac{t_j}{1+t_j}\to0,
   \quad \|\operatorname{Ric}_{g_j}+g_j\|_2
      =2t_j\sqrt{N_j\operatorname{Vol}(h_1)}\to\infty.        \tag{22}
  \] These manifolds already have Einstein metrics; (22) disproves an
  implication between norms, not an existence assertion.
\item
  \textbf{A local regularity theorem does not eliminate global kernels.}
  On \(\Sigma_g\times\Sigma_h\) with the standard negative Einstein
  product, a Teichmüller deformation of one factor, in
  transverse-traceless gauge, gives a nonzero tensor in \(\ker L_g\).
  Thus Theorem 2 deliberately does not replace the previous compact KE
  realization theorem: that theorem can accommodate such moduli without
  inverting a full operator. Deleting these modes requires an
  obstruction equation and an integrability argument, not merely a
  choice of notation.
\item
  \textbf{Neither gauge saturation nor a Weyl estimate is needed in
  Theorem 2.} The stronger topological input now handles curvature
  concentration. It does not supply unit-ball noncollapse, the linear
  gap, or the small unnormalized residual. The previous actual-form
  estimates can still provide geometric structure, but they do not
  replace these hypotheses.
\item
  \textbf{This is not a replacement of Jäckel's theorem.} His Theorem
  1.1 has strict negative sectional curvature, an exactly hyperbolic
  thin part, bounded Ricci derivative and small unnormalized residual;
  it allows arbitrarily small injectivity radius. Here local noncollapse
  and a topological regularity mechanism replace that thin-part
  analysis, while the global linear gap is explicit rather than obtained
  from negative sectional curvature. Neither statement subsumes the
  other.
\end{enumerate}

\hypertarget{relation-to-the-first-scientific-line}{%
\subsection{5. Relation to the first scientific
line}\label{relation-to-the-first-scientific-line}}

The previous compact realization mechanism was \[
 \text{gauge saturation plus bounded diameter}
 \Longrightarrow\text{smooth compact limit}
 \Longrightarrow\text{same-type KE existence}.
\] The present sufficient mechanism is different: \[
 \text{vanishing }\pi_2\text{ plus local noncollapse and bounded Ricci}
 \Longrightarrow\text{uniform local geometry},
\] followed by a separately verified global inverse and small mixed
residual, which imply a genuine solution on the original manifold. This
obtains Einstein existence without diameter, but not from unrestricted
Yau26 input. It identifies two distinct missing pieces in a large-scale
degeneration: local concentration control and global low-mode
solvability. Replacing diameter by either one alone would omit the
other.

Theorem 1 is the local scope extension obtained from the stronger
classical ALE exclusion. Theorem 2 and Corollary 3 are its explicit
composition with a classical solver, not a newly invented
inverse-function method or a claim of global priority. They supply a
usable sufficient theorem with all additional input stated, rather than
assigning another round of work to an unspecified global matching
problem.

\end{document}
